\documentclass[a4paper,11pt]{article}

\usepackage[a4paper,margin=2.5cm]{geometry}
\usepackage{graphicx}
\usepackage{amsmath,amssymb}
\usepackage{hyperref}
\usepackage{float}
\usepackage{multirow}
\usepackage[table]{xcolor}
\usepackage{booktabs}
\usepackage{tabularx}
\usepackage{array}
\usepackage{longtable}
\usepackage{enumitem}
\usepackage{titlesec}

\hypersetup{
    colorlinks=true,
    linkcolor=blue,
    urlcolor=blue,
    citecolor=blue,
    pdftitle={Surprise in Active Inference: Causes and Modulators}
}

\titleformat{\section}{\large\bfseries}{\thesection}{0.6em}{}
\titleformat{\subsection}{\normalsize\bfseries}{\thesubsection}{0.6em}{}

\setlength{\parindent}{0pt}
\setlength{\parskip}{0.6em}
\renewcommand{\arraystretch}{1.25}

\title{Surprise in Active Inference: Causes and Modulators}
\author{}
\date{}

\begin{document}
    \maketitle
    
    
    \section{Executive Summary}
        In active inference, surprise is not a casual psychological notion; it is the negative log probability of an outcome under the agent's generative model. An observation becomes surprising when the model assigns it low probability. Because exact surprise is usually intractable, agents minimize variational free energy, which serves as a computable upper bound on surprise.
        
        Surprise can increase when sensory input deviates from top-down predictions, when priors are too rigid or poorly calibrated, when precision weighting amplifies prediction errors, or when the model is misspecified relative to the environment. Friston's framework explains surprise reduction mainly through two coupled routes: changing beliefs and predictions through perception or learning, and changing sensory input through action.
    
    
    \section{Background: Surprise and Free Energy in Active Inference}
        In Friston's formulation, surprise is defined as
        \begin{equation}
            I(o) = -\ln p(o),
        \end{equation}
        where $o$ is an observation and $p(o)$ is the probability assigned to it by the generative model. If the model expects an outcome, surprise is low; if the outcome is improbable, surprise is high.
        
        Because the exact marginal probability $p(o)$ is often difficult to compute, active inference works with variational free energy,
        \begin{equation}
            F = D_{KL}\bigl(q(s) \parallel p(s \mid o)\bigr) - \ln p(o),
        \end{equation}
        which can also be written in accuracy-complexity form. By minimizing $F$, the agent indirectly minimizes surprise. This is the basis for perception, action, and learning in the free-energy principle.
    
    
    \section{Factors Increasing Surprise}
        Several factors increase surprise under active inference:
        \begin{itemize}[leftmargin=1.8em]
	\item \textbf{Prediction mismatch:} sensory input differs from the predicted input generated by higher-level beliefs.
    \item \textbf{Rigid or poor priors:} the model strongly expects one class of states and cannot accommodate alternatives.
    \item \textbf{Incorrect precision weighting:} prediction errors are overweighted or underweighted.
    \item \textbf{Model misspecification:} the generative model lacks the latent causes needed to explain the observation.
    \item \textbf{Novel structure in the environment:} the world has changed in a way that learned regularities no longer capture.
        \end{itemize}
        
        \begin{table}[H]
            \centering
            \small
            \begin{tabularx}{\textwidth}{>{\raggedright\arraybackslash}p{3.2cm} >{\raggedright\arraybackslash}X >{\raggedright\arraybackslash}X}
		\toprule
        \rowcolor{gray!15}
        \textbf{Factor} & \textbf{Mathematical role} & \textbf{Effect on surprise} \\
        \midrule
        Prediction mismatch & Large sensory prediction error under the likelihood model & Raises free energy and signals poor explanatory fit \\
        Rigid prior beliefs & Strong prior concentration on a narrow state space & Makes deviations appear highly improbable \\
        Misweighted precision & Changes the gain on prediction errors & Can exaggerate or suppress the impact of mismatch \\
        Model misspecification & Missing hidden causes or wrong dependencies & Prevents accurate inference even with updating \\
        Environmental novelty & True change in external structure or contingencies & Produces sustained surprise until learning occurs \\
        \bottomrule
            \end{tabularx}
            \caption{Main factors that increase surprise in active inference.}
        \end{table}
    
    
    \section{Likelihood, Prior, Posterior, and Precision}
        Surprise depends on the relation between four core elements.
        
        \subsection{Likelihood}
            The likelihood $p(o \mid s)$ specifies how hidden states $s$ generate observations $o$. If an observation is not well explained by likely states, surprise rises.
        
        \subsection{Prior}
            The prior $p(s)$ expresses what the agent expects before seeing the current data. Strong priors can stabilize inference, but if they are inaccurate they can also magnify surprise.
        
        \subsection{Posterior}
            The posterior $p(s \mid o)$ updates beliefs after data arrive. In variational inference, the agent uses an approximate posterior $q(s)$ and minimizes divergence between $q(s)$ and the true posterior.
        
        \subsection{Precision}
            Precision is inverse variance. In predictive coding accounts of active inference, precision weights prediction errors and controls their effective influence on belief updating. Precision therefore modulates not the raw mismatch itself, but the importance assigned to that mismatch.
    
    
    \section{Hierarchical Models and Precision-Weighted Prediction Errors}
        Friston's account is hierarchical. Higher levels encode more abstract or slower-changing causes, while lower levels encode faster or more immediate sensory causes. Higher levels send predictions downward; lower levels return prediction errors upward.
        
        A simplified relation is
        \begin{equation}
            \varepsilon = o - \hat{o},
        \end{equation}
        where $o$ is observed input and $\hat{o}$ is predicted input. Precision-weighted prediction error is then
        \begin{equation}
            \xi = \Pi \, \varepsilon,
        \end{equation}
        with $\Pi$ denoting precision. Large $\xi$ means the mismatch is both substantial and treated as reliable.
        
        \begin{center}
            \fbox{%
                \parbox{0.9\textwidth}{%
                    \textbf{Hierarchical message flow}\newline
                    Higher-level beliefs generate top-down predictions of lower-level states or sensory causes.\newline
                    Lower levels compare those predictions with actual input and compute prediction errors.\newline
                    Prediction errors are weighted by precision and sent upward to update beliefs.}}
        \end{center}
    
    
    \section{Example Calculations of Surprise}
        If an outcome has probability $p(o)=0.5$, then
        \begin{equation}
            I(o) = -\ln(0.5) \approx 0.693.
        \end{equation}
        If instead $p(o)=0.1$, then
        \begin{equation}
            I(o) = -\ln(0.1) \approx 2.303.
        \end{equation}
        Thus, less probable observations are more surprising.
        
        \begin{table}[H]
            \centering
            \small
            \begin{tabular}{>{\centering\arraybackslash}p{3.0cm} >{\centering\arraybackslash}p{3.0cm} >{\centering\arraybackslash}p{3.0cm}}
                \toprule
                \rowcolor{gray!15}
                \textbf{Observation probability $p(o)$} & \textbf{Surprise $-\ln p(o)$} & \textbf{Interpretation} \\
                \midrule
                0.50                                    & 0.693                         & Mildly surprising       \\
                0.10                                    & 2.303                         & Clearly surprising      \\
                0.01                                    & 4.605                         & Strongly surprising     \\
                \bottomrule
            \end{tabular}
            \caption{Simple numerical examples of surprise.}
        \end{table}
    
    
    \section{Empirical and Simulation Evidence}
        Predictive coding and active inference models explain mismatch responses, repetition suppression, perceptual updating, and several forms of action selection. When the incoming signal violates a learned regularity, prediction error rises and drives belief revision or action.
        
        Simulation work in Friston's papers also shows that agents can reduce expected surprise either by updating hidden-state beliefs or by acting so that sampled inputs better fit prior expectations. In later formulations, expected free energy extends this logic to future-oriented choice under uncertainty.
    
    
    \section{Implications for Perception, Learning, and Psychopathology}
        \textbf{Perception} reduces surprise by updating current beliefs about hidden causes.\newline
        \textbf{Learning} reduces future surprise by changing model parameters and priors.\newline
        \textbf{Action} reduces surprise by changing what the agent samples from the world.\newline
        \textbf{Psychopathology} can be interpreted as maladaptive inference, for example when priors are too strong, precision is misallocated, or the generative model cannot flexibly explain input.
    
    
    \section{Comparison of Factors Affecting Surprise}
        \begin{longtable}{>{\raggedright\arraybackslash}p{3.0cm} >{\raggedright\arraybackslash}p{5.0cm} >{\raggedright\arraybackslash}p{5.2cm}}
            \toprule
            \rowcolor{gray!15}
            \textbf{Component} & \textbf{What changes} & \textbf{Result for surprise} \\
            \midrule
            \endfirsthead
            \toprule
            \rowcolor{gray!15}
            \textbf{Component} & \textbf{What changes} & \textbf{Result for surprise} \\
            \midrule
            \endhead
            Perception & Current beliefs about hidden causes & Lowers surprise by improving posterior fit to present input \\
            Learning & Parameters, structure, or priors of the model & Lowers future surprise by improving the model itself \\
            Action & Sampled sensory states or environmental relation & Lowers surprise by making input conform better to predictions \\
            Attention / precision control & Gain on prediction error signals & Changes which mismatches dominate updating \\
            Novelty & Environmental contingencies or hidden causes & Raises surprise until accommodated by inference or learning \\
            \bottomrule
        \end{longtable}
    
    
    \section{How Surprise Arises and Is Resolved}
        Surprise arises when the agent encounters observations that its current model treats as improbable. This usually appears locally as prediction error and globally as elevated free energy.
        
        In Friston's framework, surprise is mainly resolved in two coupled ways. First, the agent can change its beliefs, predictions, and model parameters through inference and learning. Second, it can change the sensory data it samples by acting on the world. Precision control, attention, and hierarchical message passing are not separate fundamental routes; they are mechanisms that shape how the two main routes operate.
    
    
    \section{References}
        \begin{itemize}[leftmargin=1.8em]
	\item Friston, K. J. (2010). The free-energy principle: a unified brain theory? \textit{Nature Reviews Neuroscience}, 11(2), 127--138.
    \item Friston, K., Kilner, J., and Harrison, L. (2006). A free energy principle for the brain. \textit{Journal of Physiology - Paris}, 100(1--3), 70--87.
    \item Feldman, H., and Friston, K. (2010). Attention, uncertainty, and free-energy. \textit{Frontiers in Human Neuroscience}, 4, 215.
    \item Friston, K., Daunizeau, J., and Kiebel, S. (2009). Reinforcement learning or active inference? \textit{PLoS ONE}, 4(7), e6421.
    \item Parr, T., and Friston, K. J. (2017). Working memory, attention, and salience in active inference. \textit{Scientific Reports}, 7, 14678.
        \end{itemize}

\end{document}